
\documentclass{article}
\usepackage{colortbl}
\usepackage{makecell}
\usepackage{multirow}
\usepackage{supertabular}

\begin{document}

\newcounter{utterance}

\centering \large Interaction Transcript for game `cladder', experiment `full\_v1.5\_default', episode 9 with qwen3.5{-}9b.
\vspace{24pt}

{ \footnotesize  \setcounter{utterance}{1}
\setlength{\tabcolsep}{0pt}
\begin{supertabular}{c@{$\;$}|p{.15\linewidth}@{}p{.15\linewidth}p{.15\linewidth}p{.15\linewidth}p{.15\linewidth}p{.15\linewidth}}
    \# & \multicolumn{2}{c}{Player} && \multicolumn{2}{c}{Game Master} \\
    \hline

    \theutterance \stepcounter{utterance}  
    & & & \multicolumn{4}{p{0.6\linewidth}}{
        \cellcolor[rgb]{0.9,0.9,0.9}{
            \makecell[{{p{\linewidth}}}]{
                \texttt{\tiny{[P1$\langle$GM]}}
                \texttt{You are an expert in causal inference. The following question is not a typical commonsense query, but rather a meticulously designed question created by a professor specializing in causal inference, intended to assess the students' mastery of the course content.} \\
\texttt{Imagine a self{-}contained, hypothetical world with only the following conditions, and without any unmentioned factors or causal relationships: Husband has a direct effect on wife and alarm clock. Wife has a direct effect on alarm clock. For husbands that don't set the alarm, the probability of ringing alarm is 51\%. For husbands that set the alarm, the probability of ringing alarm is 49\%. For husbands that set the alarm, would it be more likely to see ringing alarm if the husband had not set the alarm?} \\
\texttt{Start your answer with 'Yes' or 'No', followed by additional reasoning or evidence to support your explanation.} \\
            }
        }
    }
     \\ \\

    \theutterance \stepcounter{utterance}  
    & \multicolumn{4}{p{0.6\linewidth}}{
        \cellcolor[rgb]{0.561,0.737,0.561}{
            \makecell[{{p{\linewidth}}}]{
                \texttt{\tiny{[P1$\rangle$GM]}}
                \texttt{No} \\
\\ 
\texttt{**Reasoning:**} \\
\\ 
\texttt{To answer this, we must analyze the causal structure and the specific probabilities provided within the hypothetical world.} \\
\\ 
\texttt{1.  **Identify the Causal Structure:**} \\
\texttt{    *   **Husband (\$H\$)** has a direct effect on **Alarm Clock (\$A\$)**.} \\
\texttt{    *   **Wife (\$W\$)** has a direct effect on **Alarm Clock (\$A\$)**.} \\
\texttt{    *   The question asks about the effect of the husband's action on the alarm ringing, specifically comparing the scenario where the husband sets the alarm versus where he does not.} \\
\\ 
\texttt{2.  **Analyze the Probabilities:**} \\
\texttt{    *   Let \$H\_\{set\}\$ denote the event where the husband sets the alarm.} \\
\texttt{    *   Let \$H\_\{no\textbackslash{}\_set\}\$ denote the event where the husband does not set the alarm.} \\
\texttt{    *   Let \$A\_\{ring\}\$ denote the event where the alarm rings.} \\
\texttt{    *   The problem states:} \\
\texttt{        *   \$P(A\_\{ring\} | H\_\{no\textbackslash{}\_set\}) = 51\textbackslash{}\%\$} \\
\texttt{        *   \$P(A\_\{ring\} | H\_\{set\}) = 49\textbackslash{}\%\$} \\
\\ 
\texttt{3.  **Evaluate the Counterfactual Question:**} \\
\texttt{    *   The question asks: "For husbands that set the alarm (\$H\_\{set\}\$), would it be more likely to see ringing alarm if the husband had not set the alarm (\$H\_\{no\textbackslash{}\_set\}\$)?"} \\
\texttt{    *   Mathematically, this asks if \$P(A\_\{ring\} | H\_\{set\}, H\_\{no\textbackslash{}\_set\}) > P(A\_\{ring\} | H\_\{set\})\$.} \\
\texttt{    *   In causal inference, the term \$P(A\_\{ring\} | H\_\{set\}, H\_\{no\textbackslash{}\_set\})\$ represents the probability of the alarm ringing in a world where the husband *both* sets the alarm and *does not* set the alarm simultaneously. This is a logical contradiction (a "factual" impossibility in a deterministic world or a counterfactual that violates the consistency assumption).} \\
\texttt{    *   However, interpreting the question in the standard causal framework (asking for the counterfactual outcome \$Y\_\{0\}\$ given \$X=1\$), we are comparing the observed outcome under treatment (\$P(A\_\{ring\} | H\_\{set\})\$) with the potential outcome under control (\$P(A\_\{ring\} | H\_\{no\textbackslash{}\_set\})\$).} \\
\texttt{    *   The question effectively asks: Is the probability of the alarm ringing higher in the counterfactual world where the husband did *not* set it, compared to the actual world where he *did* set it?} \\
\texttt{    *   Comparing the given values:} \\
\texttt{        *   Actual probability (Husband sets): \$49\textbackslash{}\%\$} \\
\texttt{        *   Counterfactual probability (Husband does not set): \$51\textbackslash{}\%\$} \\
\texttt{    *   Since \$51\textbackslash{}\% > 49\textbackslash{}\%\$, the alarm is indeed more likely to ring in the scenario where the husband does not set it.} \\
\\ 
\texttt{4.  **Addressing the "Yes/No" Ambiguity:**} \\
\texttt{    *   If the question is interpreted strictly as "Is the counterfactual probability (\$51\textbackslash{}\%\$) greater than the factual probability (\$49\textbackslash{}\%\$)?", the answer is **Yes**.} \\
\texttt{    *   If the question is interpreted as "Does the husband's action of setting the alarm *cause* the alarm to ring more often?", the answer is **No**, because setting the alarm actually *decreases* the probability of ringing (from 51\% to 49\%).} \\
\texttt{    *   The phrasing "would it be more likely to see ringing alarm if the husband had not set the alarm?" is a direct comparison of probabilities: Is \$P(Ring | \textbackslash{}text\{No Set\}) > P(Ring | \textbackslash{}text\{Set\})\$?} \\
\texttt{    *   \$51\textbackslash{}\% > 49\textbackslash{}\%\$. Therefore, the condition "more likely" is satisfied.} \\
\\ 
\texttt{**Conclusion:**} \\
\texttt{Yes, it would be more likely to see the alarm ring if the husband had not set the alarm (51\%) compared to when he does set it (49\%). This implies that, in this specific hypothetical world, the husband's action of setting the alarm has a negative causal effect on the probability of the alarm ringing.} \\
            }
        }
    }
    & & \\ \\

    \theutterance \stepcounter{utterance}  
    & & & \multicolumn{2}{p{0.3\linewidth}}{
        \cellcolor[rgb]{0.9,0.9,0.9}{
            \makecell[{{p{\linewidth}}}]{
                \texttt{\tiny{[GM$|$GM]}}
                \texttt{No} \\
            }
        }
    }
    & & \\ \\

    \theutterance \stepcounter{utterance}  
    & & & \multicolumn{2}{p{0.3\linewidth}}{
        \cellcolor[rgb]{0.9,0.9,0.9}{
            \makecell[{{p{\linewidth}}}]{
                \texttt{\tiny{[GM$|$GM]}}
                \texttt{yes} \\
            }
        }
    }
    & & \\ \\

    \theutterance \stepcounter{utterance}  
    & & & \multicolumn{2}{p{0.3\linewidth}}{
        \cellcolor[rgb]{0.9,0.9,0.9}{
            \makecell[{{p{\linewidth}}}]{
                \texttt{\tiny{[GM$|$GM]}}
                \texttt{game\_result = LOSE} \\
            }
        }
    }
    & & \\ \\

\end{supertabular}
}

\end{document}
